\section{结论与启示}

本文围绕“数字化转型还是叙事包装”这一问题，提出了一套基于 IPO 招股说明书与问询回复文本识别企业 AI 披露真实性的研究框架。与仅统计关键词频次或只做案例叙述的做法不同，本文将提及强度、场景具体性和表述审慎性结合起来构建真实性综合指数，并同时保留包装性这一反向面向，形成“真实性-包装性”双面识别结构。方法上，本文采用关键词筛选、相关段落抽取、LLM 辅助标注、人工复核、指数构建和基础回归验证的轻量化路径，力图在本科阶段能力边界内实现研究规范性与执行可行性的平衡。

本文预期的核心结论是：若企业在 IPO 阶段对 AI 与数字化的披露更具体、更审慎、且能在问询回复中提供可验证证据，则其披露更可能体现为真实能力信号，并在上市后经营结果中得到兑现；相反，若企业披露主要表现为抽象战略口号、宏大叙事和承诺化措辞，则这类表述更接近叙事包装，其对上市后经营改善的支持力度较弱。无论最终回归结果的显著性强弱如何，本文至少能够帮助区分“说得多”和“说得实”之间的差别，从而把企业数字化叙事研究从简单词频统计推进到更具问题意识的真实性识别。

本文可能的启示主要体现在三个方面。对投资者而言，识别企业 AI 叙事时不应只关注是否频繁提及技术热词，更应关注其是否落到具体业务场景、组织投入和实施边界。对监管部门而言，问询环节中围绕数字化能力真实性、先进性和必要性的追问，具有重要的信息治理价值。对企业而言，数字化叙事若要形成可信信号，需要依托真实投入、系统建设和组织能力，而非仅依赖概念包装。

本文也存在若干局限。首先，真实性和包装性的识别仍然依赖一定程度的人工判断，尽管本文已通过 LLM 辅助和人工复核降低主观性，但评分误差仍难完全避免。其次，上市后经营表现受到多种因素影响，本文采用的是相对简洁的关联性识别框架，并不直接主张严格因果解释。最后，受数据可得性限制，正式实证的时间窗口和样本规模仍需在执行过程中根据数据质量进行平衡。未来研究可以在扩大样本、引入更多客观技术投入指标、延长经营结果观察期和优化半自动文本识别规则等方面继续推进。
